% frontmatter/abstract-en.tex — Abstract (Civil Engineering perspective)
\chapter*{Abstract}
\addcontentsline{toc}{chapter}{Abstract}

\noindent\textbf{Title:} \thesistitleEN\\
\textbf{By:} \authornameEN\\
\textbf{Advisor:} \advisorEN, Ph.D.\\

Construction cost estimation is a critical activity in project management,
particularly for public works in which quantity surveyors must reference the
official median prices issued by the Comptroller General's Department (CGD)
and the construction materials price index of the Ministry of Commerce in
order to prepare the Bill of Quantities (BOQ) used for tendering, procurement,
and project cost control. In current practice, however, three persistent
problems remain: (1) material price information is dispersed across many
unconnected sources, (2) government-issued reference prices are updated on
a monthly cycle that lags behind the actual market, and (3) the published
prices are typically national averages that do not reflect the considerable
spatial variability between provinces.

This project develops an information system for material cost estimation
that integrates price data from eleven Thai sources --- three public
agencies (the Trade Policy and Strategy Office, the Comptroller General's
Department, and the Department of Internal Trade) and eight modern-trade
retailers (HomePro, MegaHome, Boonthavorn, Thai Watsadu, Global House,
Dohome, SCG Home, and BnB Home). The system computes unit costs in
accordance with the CGD median-price methodology for five common
construction work types: wall tiling, lintel/post forming, reinforcing
steel, painting, and brickwork. From a user-supplied work quantity, the
system generates a detailed Bill of Materials and a total cost breakdown
suitable for use in BOQ documentation.

For data quality, the system performs data cleaning, unit standardization,
cross-source item matching using canonical specifications (brand, size,
grade), and statistical outlier detection based on z-score testing
following the approach of Lee and Yun (2024), so that prices distorted
by technical errors or short-term promotions are excluded from the cost
calculation.

Comparative analysis based on real data collected by the system
(snapshot: 24 May 2026, Bangkok) shows that DIT prices exceed the CGD
baseline by a median of $+$8.4\% (mean $+$10.1\%) across sixteen comparable
price pairs (same unit, same product category), that prices in the Southern region
are 22.9\% higher than those in the Central region (175 vs 215 baht per
50 kg bag of Portland cement Type I), consistent with long-distance
transportation costs, and that the TPSO price series for cement in
Bangkok shows a linear upward trend of 0.819 baht per data point over
ten months ($R^{2} = 0.871$). Price data from modern-trade retailers
is currently being collected through an AI agent research pipeline
under active development.

The resulting system serves as a working tool for civil engineers,
quantity surveyors, architects, and public officials, allowing
multi-source price comparison within a single application. It reduces
the time required for price research, improves the accuracy of cost
estimates, and supports greater transparency in the public procurement
process by exposing comparable price information from independent
sources.

\noindent\textbf{Keywords:} Construction Cost Estimation,
Material Unit Cost, Bill of Quantities (BOQ), Median Price (CGD),
Construction Materials Price Index, Modern Trade, Item Canonicalization,
Outlier Detection
\newpage
